% ---------------------------------------------------------------------------
% Preambulo de ajuste para IEEEtran sobre Quarto/Pandoc
% ---------------------------------------------------------------------------
% Se incluye al final del preambulo (include-in-header), despues de que
% Pandoc haya cargado sus paquetes por defecto.

% 1) Pandoc carga el paquete `parskip` en clases no-KOMA, lo que anula la
%    sangria de parrafo. IEEEtran define su propio espaciado, asi que se
%    restauran los valores tipograficos de IEEE (parrafo con sangria, sin
%    espacio adicional entre parrafos).
\setlength{\parindent}{1em}
\setlength{\parskip}{0pt}
\setlength{\parsep}{0pt}

% 2) El paquete `caption` se carga DESPUES de este preambulo, por lo que el
%    ajuste se difiere a \AtBeginDocument. Desconoce IEEEtran y aplica su
%    estilo por defecto ("Figura 1: ..."); se aproxima la convencion IEEE:
%    etiqueta en negrita y separada por punto.
\AtBeginDocument{\captionsetup{labelfont=bf,labelsep=period,font=small}}
\AtBeginDocument{\captionsetup[figure]{font=footnotesize}}

% 3) Quarto/Pandoc emite \author{...} solo con el nombre. Se reemplaza por el
%    bloque de autoria de IEEEtran (nombre, afiliacion y correo) justo antes de
%    que \maketitle lo consuma.
\makeatletter
\AtBeginDocument{%
  \def\@IEEEauthorblockAinterlinespace{2.5ex}%
  \author{%
    \IEEEauthorblockN{Deisy Elizabeth Calle Calle}%
    \IEEEauthorblockA{Carrera de Ingeniería en Telemática, Facultad de Informática y Electrónica,\\
    Escuela Superior Politécnica de Chimborazo (ESPOCH), Riobamba, Ecuador\\
    calleelizabeth640@gmail.com\\[-4pt]
    {\footnotesize Código 36 --- Sistemas Distribuidos y Paradigmas de la Computación (ETL1TP50)}}%
  }%
  \def\IEEEkeywordsname{Palabras clave}%
}
\makeatother

% 4) Ajuste fino del espaciado alrededor de los flotantes para que cuatro
%    figuras y dos tablas convivan con el texto sin paginas casi vacias.
\setlength{\textfloatsep}{8pt plus 2pt minus 2pt}
\setlength{\intextsep}{8pt plus 2pt minus 2pt}
\setlength{\floatsep}{8pt plus 2pt minus 2pt}
